% =============================================================================
%   ABSTRACT
% =============================================================================
\begin{abstract}
Mid-cap investors increasingly demand ESG-aware allocation tools, yet
practical index methods remain constrained by sparse and inconsistent
non-financial disclosure outside large-cap universes. This study develops a
transparent 9-factor ESG-integrated composite index for 276 mid-cap firms
across US and Indian markets using hybrid ESG inputs that combine real Yahoo
Finance governance risk scores and SEC EDGAR governance disclosures with
proxy-derived environmental and social indicators. The balanced multi-factor
specification ranks higher than single-factor alternatives on cross-sectional
quality proxies and demonstrates superior cross-sectional quality in stressed
market decompositions and rolling-rebalance comparisons. ESG integration
functions primarily as a risk filter within the broader multi-factor system,
with higher ESG associated with lower volatility and beta rather than a
standalone alpha effect. Robustness checks indicate that the ranking framework
remains stable under cross-validation at the portfolio level (individual rank confidence intervals remain wide, spanning 247--255 positions, indicating that rank stability applies to aggregate portfolio characteristics rather than individual company positions), and high-cap generalization testing on 45 S\&P 500 companies confirms full transferability (Kendall $W = 0.933$, 9/9 factors pass z-score test, Spearman $\rho = 0.9993$). Monotonicity evidence shows 7 of 8 factors exhibit monotonic quintile returns against the forward quality proxy, while 2 of 9 achieve statistical significance against the trailing return proxy (Jonckheere-Terpstra test, $p < 0.05$). All
performance metrics reflect cross-sectional factor quality, not realized
portfolio returns.
\end{abstract}

\begin{IEEEkeywords}
ESG integration, multi-factor model, mid-cap equities, composite index,
sustainability investing
\end{IEEEkeywords}
